% Nuestra definición


En este trabajo se van a definir los transformers como ¿máquinas aceptadoras?. \todo{¿Cuál sería la palabra correcta?}
Un transformer consta de una cinta de $n_{max}$ celdas de la que lee el input (escrito en las primeras $n$ celdas) y va escribiendo a derecha de lo ya escrito nuevos caracteres usando un generador de tokens.

El generador de tokens está compuesto por 4 partes: una función de token embedding (TE), una función de codificación posicional (PE), $L$ capas de decodificación y una capa final de output.
Se define $L$ como la profundidad del transformer.
Cada capa de decodificación está compuesta por dos subcapas: una capa de atención (ATTN) y una red de propagación hacia adelante totalmente conectada por posiciones (\textit{a position-wise fully-connected feed-forward network}) (FF).

Cada transformer opera sobre un alfabeto $\Sigma$ y está definido por los parámetros $\theta = (\theta_{TE}, \theta_{PE}, \{\theta_{ATTN}^{(l)}, \theta_{FF}^{(l)}\}_{l=0}^{L-1}, \theta_{output})$. \todo{¿Un comentario sobre que $\theta$ es lo que se entrena?}

El modelo estudiado en este trabajo es no uniforme, es decir que para palabras de distintas longitudes se utilizan distintos transformers.

% Definición de cada parte:
Token embedding: es un mapeo de los elementos de $\Sigma$ a vectores de $\mathbb{R}^d$. El mapeo está definido por $\theta_{TE} \in \mathbb{R}^{d \times \mid \Sigma \mid}$. Para todo $a \in \Sigma$ se va a notar $\theta_{TE}(a) \in \mathbb{R}^d$ al mapeo del caracter $a$.

Codificación posicional: es un mapeo de las posiciones de la cinta a vectores de $\mathbb{R}^d$. El mapeo está definido por $\theta_{TE} \in \mathbb{R}^{d \times n_{max}}$. Para todo $1 < i \leq n_{max}$ se va a notar $\theta_{PE}(i) \in \mathbb{R}^d$ al mapeo de la posición $i$.

Capa de atención: dados los parámetros $\theta_{ATTN} = (W_Q, W_K , W_V , W_O) \in \mathbb{R}^d \times \mathbb{R}^d \times \mathbb{R}^d \times \mathbb{R}^d$ la capa de atención se define como el mapeo definido en el algoritmo ?? de $\mathbb{R}^{d} \times \dots \mathbb{R}^d$ a $\mathbb{R}^{d} \times \dots \mathbb{R}^d$. \todo{¿Cómo noto que son n vectores de $R^d$?}

[Acá viene el algoritmo 1 de toyota \todo{ver cómo copiarlo}]

Capa de feed forward?\todo{Hay que ver cómo nombrar esta capa}: dado $\theta_{FF} = (W_1 , b_1 , W_2 , b_2) \in \mathbb{R}^{d \times d} \times \mathbb{R}^d \times \mathbb{R}^{d \times d} \times \mathbb{R}^d$, se define la capa de feed forward como un mapeo de $\mathbb{R}^d$ a $\mathbb{R}^d$ dado por la siguiente ecuación $FF_{\theta_{FF}} := W_2 \; \relu(W_1 h + b1 ) + b2$. Donde $\relu: \mathbb{R}^d \to \mathbb{R}^d$ reemplazando cada coordenada por el máximo entre su valor y $0$: $relu(h)_i = max(h_i, 0)$ para todo $0 < i \leq d$.

Output: dado $\theta_{output} \in \mathbb{R}^{\mid \Sigma \mid \times d}$ la capa de output se define como un mapeo de $\mathbb{R}^d$ a una distribución de probabilidad sobre $\Sigma$. Para todo $h \in \mathbb{R}^d$ el mapeo se computa como $\softmax(\theta_{output}h)$.



% Algoritmo con la estructura del generador


% Definición de CoT


% Definición de computar un lenguaje


% Comentarios sobre la definición de toyota y bisimulación

    % Diferencias:
    %   - Nosotros no tenemos el paso de output porque es al pedo
    %   - 
